\documentclass[10pt,a4paper]{amsart}
\usepackage{amsmath,amssymb,amsthm,graphicx}
% Retained preamble for the Stacks Project chapter snapshot in this directory.
% The chapter's own preamble.tex is a copy of the shared Stacks Project preamble; the block
% below is sliced MECHANICALLY, line for line, from it. Nothing was added, removed or
% redefined; the wrapper class, the xr-hyper/hyperref wiring and the externaldocument list
% are omitted because the assembler supplies the document class.
%
% lines taken from preamble.tex: 3, 5, 9, 11, 12, 156, 157, 158, ... (51 lines total)
%
% package coverage restored from the chapter's own preamble (2026-10-13): xy, xy-2cell

\usepackage[all]{xy}
\xyoption{2cell}
\UseAllTwocells

\newenvironment{reference}{\comment}{\endcomment}
\newenvironment{slogan}{\comment}{\endcomment}
\newenvironment{history}{\comment}{\endcomment}
\theoremstyle{plain}
\newtheorem{theorem}[subsection]{Theorem}
\newtheorem{proposition}[subsection]{Proposition}
\newtheorem{lemma}[subsection]{Lemma}
\theoremstyle{definition}
\newtheorem{definition}[subsection]{Definition}
\newtheorem{example}[subsection]{Example}
\newtheorem{exercise}[subsection]{Exercise}
\newtheorem{situation}[subsection]{Situation}
\theoremstyle{remark}
\newtheorem{remark}[subsection]{Remark}
\newtheorem{remarks}[subsection]{Remarks}
\numberwithin{equation}{subsection}
\def\lim{\mathop{\mathrm{lim}}\nolimits}
\def\colim{\mathop{\mathrm{colim}}\nolimits}
\def\Spec{\mathop{\mathrm{Spec}}}
\def\Hom{\mathop{\mathrm{Hom}}\nolimits}
\def\Ext{\mathop{\mathrm{Ext}}\nolimits}
\def\SheafHom{\mathop{\mathcal{H}\!\mathit{om}}\nolimits}
\def\SheafExt{\mathop{\mathcal{E}\!\mathit{xt}}\nolimits}
\def\Sch{\mathit{Sch}}
\def\Mor{\mathop{\mathrm{Mor}}\nolimits}
\def\Ob{\mathop{\mathrm{Ob}}\nolimits}
\def\Sh{\mathop{\mathit{Sh}}\nolimits}
\def\NL{\mathop{N\!L}\nolimits}
\def\CH{\mathop{\mathrm{CH}}\nolimits}
\def\proetale{{pro\text{-}\acute{e}tale}}
\def\etale{{\acute{e}tale}}
\def\QCoh{\mathit{QCoh}}
\def\Ker{\mathop{\mathrm{Ker}}}
\def\Im{\mathop{\mathrm{Im}}}
\def\Coker{\mathop{\mathrm{Coker}}}
\def\Coim{\mathop{\mathrm{Coim}}}
\DeclareMathSymbol{\boxtimes}{\mathbin}{AMSa}{"02}
\def\QCohstack{\mathcal{QC}\!\mathit{oh}}
\def\Cohstack{\mathcal{C}\!\mathit{oh}}
\def\Spacesstack{\mathcal{S}\!\mathit{paces}}
\def\Quotfunctor{\mathrm{Quot}}
\def\Hilbfunctor{\mathrm{Hilb}}
\def\Curvesstack{\mathcal{C}\!\mathit{urves}}
\def\Polarizedstack{\mathcal{P}\!\mathit{olarized}}
\def\Complexesstack{\mathcal{C}\!\mathit{omplexes}}
\def\Pic{\mathop{\mathrm{Pic}}\nolimits}
\def\Picardstack{\mathcal{P}\!\mathit{ic}}
\def\Picardfunctor{\mathrm{Pic}}
\def\Deformationcategory{\mathcal{D}\!\mathit{ef}}
\usepackage{cleveref}
\crefname{originaldefinition}{original definition}{original definitions}
\Crefname{originaldefinition}{Original definition}{Original definitions}
\makeatletter
\newlabel{descent-definition-descending-types-morphisms}{{02W2}{source}{}{sourcepointer-1}{}}
\newlabel{descent-definition-descending-types-morphisms@cref}{{[originaldefinition][0][]02W2}{[1][0][]source}{}{}{}}
\makeatother
\crefname{originaldefinition}{original definition}{original definitions}
\Crefname{originaldefinition}{Original definition}{Original definitions}
\makeatletter
\newlabel{descent-definition-descent-datum}{{023V}{source}{}{sourcepointer-2}{}}
\newlabel{descent-definition-descent-datum@cref}{{[originaldefinition][0][]023V}{[1][0][]source}{}{}{}}
\makeatother
\crefname{originallemma}{original lemma}{original lemmas}
\Crefname{originallemma}{Original lemma}{Original lemmas}
\makeatletter
\newlabel{descent-lemma-descent-data-sheaves}{{02W5}{source}{}{sourcepointer-3}{}}
\newlabel{descent-lemma-descent-data-sheaves@cref}{{[originallemma][0][]02W5}{[1][0][]source}{}{}{}}
\makeatother
\crefname{originallemma}{original lemma}{original lemmas}
\Crefname{originallemma}{Original lemma}{Original lemmas}
\makeatletter
\newlabel{lemma-representable-morphisms-spaces-property}{{02WZ}{source}{}{sourcepointer-4}{}}
\newlabel{lemma-representable-morphisms-spaces-property@cref}{{[originallemma][0][]02WZ}{[1][0][]source}{}{}{}}
\makeatother
\crefname{originallemma}{original lemma}{original lemmas}
\Crefname{originallemma}{Original lemma}{Original lemmas}
\makeatletter
\newlabel{lemma-representable-over-space}{{02WY}{source}{}{sourcepointer-5}{}}
\newlabel{lemma-representable-over-space@cref}{{[originallemma][0][]02WY}{[1][0][]source}{}{}{}}
\makeatother
\title{A sheaf mapping representably with effective descent to an algebraic space is an algebraic space}
\author{The Stacks Project Authors}
\date{stacks-project:spaces.tex:local-snapshot-20261012; GNU FDL 1.2 or later}
\begin{document}
\maketitle
\setcounter{subsection}{1}
\noindent\textit{Editorial note.} All delivered mathematical text is the author's own, copied from the pinned Stacks Project chapter file spaces.tex; no statement or proof step was rewritten, reordered or removed.\par\medskip
\noindent\textit{Editorial note.} The proof contains a multi-line align and a bundle-like display; both are part of the counted proof span and are preserved verbatim.\par\medskip
\noindent\textit{Editorial note.} The five references of the counted spans are recorded as non-load-bearing navigation pointers; the argument that consumes them stays inside the counted proof span.\par\medskip
\noindent\textit{Editorial note.} Preamble: sources/stacks-spaces/retained-preamble.tex, sliced mechanically from the chapter's own preamble.tex.\par\medskip
\noindent\textit{Allowed standard prerequisites (editorial).}\begin{itemize}
\item Effective descent for morphisms of a type of property P. Descent, Definition descent-definition-descending-types-morphisms, used at proof lines 2010-2013 and 2051-2055. Is the hypothesis that makes the descent datum representable.
\item Descent data for sheaves are effective. Descent, Lemma descent-lemma-descent-data-sheaves, used at proof lines 2043-2050. Identifies the descent datum with the one attached to the sheaf G x\_F T.
\item A representable transformation over a fixed algebraic space has algebraic space source. Lemma lemma-representable-over-space (1926-1936), used at proof lines 2057-2059. Concludes that G is an algebraic space.
\item Properties of representable morphisms of spaces. Lemma lemma-representable-morphisms-spaces-property, used at proof lines 2059-2061. Transfers property P to G -> F.
\end{itemize}\medskip
\ifdefined\hypertarget\hypertarget{sourcepointer-1}{}\fi
\noindent\textit{Source reference.} Original definition 02W2: focused-seeds/AGP/sources/stacks-snapshot-20261011/tags.txt, source lines 3789--3789. Navigation pointer into the Stacks Project; the standard fact it supplies is applied inside the counted proof span and the argument resting on it stays there.\par\medskip
\ifdefined\hypertarget\hypertarget{sourcepointer-2}{}\fi
\noindent\textit{Source reference.} Original definition 023V: focused-seeds/AGP/sources/stacks-snapshot-20261011/tags.txt, source lines 2756--2756. Navigation pointer into the Stacks Project; the standard fact it supplies is applied inside the counted proof span and the argument resting on it stays there.\par\medskip
\ifdefined\hypertarget\hypertarget{sourcepointer-3}{}\fi
\noindent\textit{Source reference.} Original lemma 02W5: focused-seeds/AGP/sources/stacks-snapshot-20261011/tags.txt, source lines 3792--3792. Navigation pointer into the Stacks Project; the standard fact it supplies is applied inside the counted proof span and the argument resting on it stays there.\par\medskip
\ifdefined\hypertarget\hypertarget{sourcepointer-4}{}\fi
\noindent\textit{Source reference.} Original lemma 02WZ: focused-seeds/AGP/sources/stacks-snapshot-20261011/tags.txt, source lines 3822--3822. Navigation pointer into the Stacks Project; the standard fact it supplies is applied inside the counted proof span and the argument resting on it stays there.\par\medskip
\ifdefined\hypertarget\hypertarget{sourcepointer-5}{}\fi
\noindent\textit{Source reference.} Original lemma 02WY: focused-seeds/AGP/sources/stacks-snapshot-20261011/tags.txt, source lines 3821--3821. Navigation pointer into the Stacks Project; the standard fact it supplies is applied inside the counted proof span and the argument resting on it stays there.\par\medskip
\begin{lemma}
\label{lemma-morphism-sheaves-with-P-effective-descent-etale}
Let $S$ be a scheme contained in $\Sch_{fppf}$.
Let $G \to F$ be a transformation of presheaves on $(\Sch/S)_{fppf}$.
Let $\mathcal{P}$ be a property of morphisms of schemes.
Assume
\begin{enumerate}
\item $\mathcal{P}$ is preserved under any base change, fppf local on the
base, and morphisms of type $\mathcal{P}$ satisfy descent for fppf coverings,
see Descent, Definition \ref{descent-definition-descending-types-morphisms},
\item $G$ is a sheaf,
\item $F$ is an algebraic space,
\item there exists a $U \in \Ob((\Sch/S)_{fppf})$
and a surjective \'etale morphism $U \to F$ such that
$V = G \times_F U$ is representable, and
\item $V \to U$ has $\mathcal{P}$.
\end{enumerate}
Then $G$ is an algebraic space, $G \to F$ is representable and has property
$\mathcal{P}$.
\end{lemma}

\begin{proof}
Let $T$ be a scheme and let $T \to F$ be a morphism. Then
$U \times_F T \to T$ is surjective \'etale, hence $\{U \times_F T \to T\}$ is
a covering for the \'etale topology. Consider
$$
W = G \times_F (U \times_F T) = V \times_F T = V \times_U (U \times_F T).
$$
It is a scheme since $F$ is an algebraic space. The morphism
$W \to U \times_F T$ has property $\mathcal{P}$ since it is a
base change of $V \to U$. There is an isomorphism
\begin{align*}
W \times_T (U \times_F T) & =
(G \times_F (U \times_F T)) \times_T (U \times_F T) \\
& = (U \times_F T) \times_T (G \times_F (U \times_F T)) \\
& = (U \times_F T) \times_T W
\end{align*}
over $(U \times_F T) \times_T (U \times_F T)$. The middle equality maps
$((g, (u_1, t)), (u_2, t))$ to $((u_1, t), (g, (u_2, t)))$.
This defines a descent datum for $W/U \times_F T/T$, see
Descent, Definition \ref{descent-definition-descent-datum}.
This follows from
Descent, Lemma \ref{descent-lemma-descent-data-sheaves}.
Namely we have a sheaf $G \times_F T$, whose
base change to $U \times_F T$ is represented by $W$ and the isomorphism
above is the one from the proof of
Descent, Lemma \ref{descent-lemma-descent-data-sheaves}.
By assumption on $\mathcal{P}$ the descent datum above is representable.
Hence by the last statement of
Descent, Lemma \ref{descent-lemma-descent-data-sheaves}
we see that $G \times_F T$ is representable. This proves that
$G \to F$ is a representable transformation of functors.

\medskip\noindent
As $G \to F$ is representable, we see that $G$ is an algebraic space by
Lemma \ref{lemma-representable-over-space}. The fact that $G \to F$ has
property $\mathcal{P}$ now follows from
Lemma \ref{lemma-representable-morphisms-spaces-property}.
\end{proof}

\end{document}
